\documentclass[../main]{subfiles}
\begin{document}
\ifSubfilesClassLoaded{\setcounter{page}{14}}{}
\section{The Foundation of Communism}
\label{sec:03_foundation}

The foundation of capitalism is commodity production. All the wealth of capitalist society appears as an aggregate of commodities. Since the products of labour under capitalism inevitably take on the commodity form, the commodity is the “cell of capitalism”.

In accordance with the materialist conception of history, the foundation of communism as a social order can likewise only be a corresponding mode of production.

All the wealth of communist society lies in the universality of the content and forms of human activity. Marxism in general sees the essence of man in activity aimed at transforming nature (including his own social nature). To labour, to work up natural material, to create things that do not exist in nature—this is what fundamentally distinguishes human life from the life activity of an animal and determines all the other distinctions. Yet until now labour itself, on a social scale, has been unfree. Every pre-communist mode of production was marked by its own specific form of the unfreedom of labour. Communism, by contrast, presupposes the complete overcoming of this unfreedom. Therefore its foundation can only be free labour “irrespective of quotas and without reward”, as man’s “life’s prime want”.

The simplest directly collective form in which such labour is carried out (or, better said, such activity, since the word “labour” still retains a negative connotation to some extent) is the very “cell” of communism, which contains within it the further development of communist society in all its necessary determinations.

Free activity and its directly collective form are a reality in our time, just as commodity production and the commodity were a reality even before the emergence of capitalism. And the transition to a new social system will take place when man’s free activity becomes a universal relation, just as the commodity once did. From the motley array of commodities on the market, each with its own properties, one cannot directly grasp what a commodity and commodity production are. Likewise, from the diversity of types and methods of organising communist activity in their particularity, one cannot grasp what it is and what its directly collective form is. All the more so because words can obscure the essence of the matter. The directly collective form of free activity can exist even when a person carries it out entirely alone. We must single out those determinations that are necessary to it.

In the directly collective form, free activity is carried out by the measure of freedom, truth, goodness and beauty, and hence of the highest human interest; that is, it is a form of sensuous activity valuable in and of itself. Free activity necessarily presupposes the unity of all three components—sensuousness, morality and thinking—and their elevation to the level of freedom. In the course of such activity, man is affirmed in all his fullness, regardless of what exactly is being done and of which element (truth, goodness or beauty) comes to the fore.

This directly collective form is merely the activity itself, together with its corresponding content, viewed from the side of form. In everyday speech, such activity is often called social creativity, or simply creativity. Form matters because development takes place and finds expression in the necessary movement of form. Content always calls into being a form corresponding to it and is realised through that form. That is why the contradictions of content are first grasped by consciousness as contradictions of form. It is in the unfolding of form, and through that unfolding (the succession of forms), that the movement of content actually takes place. On the side of form, the important aspects of activity are the nature of the ties between people, the goal, and the social being of the object and means of labour.

Direct collectivity in free activity is expressed, first and foremost, in its orientation towards other people. In directly collective forms of activity, the human being always appears as the goal and the highest value. Everything that is done is done directly for the development of other people. For this reason, the product of such labour never acquires an independent thing-like form (the form of a commodity). The product of such labour cannot confront living people as a thing alien to them. In other words, this activity is directed towards the development of society and the improvement of other people’s lives; it is carried out for them, and therefore also for oneself as a member of that society. What are now customarily called material values can serve only as means in free communist activity.

As for the nature of the ties between people, it can, by and large, be defined as comradeship, or, what amounts to the same thing, free association.

The nature of these ties also points to the means of such activity, since in a free association none of the participants can serve as a means for carrying out someone else’s will or interests alien to him. Even if someone, thanks to his abilities and skills, takes the lead in the collective, the others do not blindly carry out another’s will; rather, they cooperate in accomplishing a common cause that is equally important and necessary for all its participants.

The nature of these ties also points to the social being of the means and subject of communist activity. Thus, only things and processes distinct from human beings can serve as means. The social form of being of the means and subjects of labour presupposes their directly social use. It was in precisely this sense that Marx spoke of the transition from the government of persons to the administration of things.

From what has been said above, it is clear that such forms of human activity are free of any external coercion whatsoever, whether direct or indirect (remuneration for labour), since neither the process nor the result of communist activity is a matter of indifference to the one who carries it out. In each instance, a person has a direct interest in his own labour, unlike the situation that is normal for modern society, in which the activity itself serves merely as a means of obtaining money.

Of course, direct collectivity also consists in the fact that the labour of any member of society rests on the labour of other people, that the subject, means and method of labour are products created by other people, while his own labour, in turn, is an integral part of all social labour and provides material for future activity. But this is characteristic of any labour carried out in any social form, not only of communist activity, and therefore it does not concern what is specific to communism.

By the directly collective form of free activity we mean the association of people in the process of activity, where the goal of each is the development of all, and the means of achieving it are things and natural processes with respect to which free and equal expression of will takes place.

The directly collective form of activity is not merely an external social form of activity, one that may or may not be present without affecting the content of the activity; on the contrary, it is its substantive form. Outside it, free activity by the standards of truth, goodness and beauty cannot be carried out at all.\footnote{What is meant here is by no means the common representation in which beauty is associated with taste, and goodness with opinions about what is good and bad, which may vary, while only truth is accorded any kind of objective significance. The categories of “truth”, “goodness” and “beauty” were developed over millennia as objective ideal characteristics of human activity in the world, unfolding (developing) within that activity, and therefore also as objective characteristics of the world (of matter). In other words, to put it briefly, when we speak of activity by the standards of truth, goodness and beauty, we are talking about the “points of growth” of human sensuousness, cognition and morality, about the development of the relationship of human to human and of human to nature, taken from the side of the ideal dimension.} The movement and metamorphoses of such a form presuppose the preservation of the aforementioned moments, just as value is always preserved in the metamorphoses of the commodity form of the products of labour (that which makes a commodity precisely a commodity, rather than simply a product of labour).

As for the content of such activity, it is always revolutionary. Herein lies its contradiction with respect to itself. Free communist activity constantly posits both its own internal limit and the transcendence of that limit. By means of historically attained, and therefore historically limited, frontiers of activity, this activity itself changes, positing both new content (thanks to the widening horizon of the assimilation of matter) and a new form. This does not happen in a roundabout way, independently of the participants in the labour process, as it does in the production of things. On the contrary, this change happens purposefully. Therefore communism, as man’s assimilation of his own activity, lays bare the internal contradictoriness of human activity in general, and hence the internal contradictoriness of the material world in which man carries out his activity. That is why Marx’s remark in the first of the Theses on Feuerbach is of paramount methodological importance for the science of communism. “The chief defect of all hitherto existing materialism—that of Feuerbach included—is that the thing, reality, sensuousness, is conceived only in the form of the object or of contemplation, but not as sensuous human activity, practice, not subjectively. Hence it happened that the active side, in contradistinction to materialism, was developed by idealism—but only abstractly, since, of course, idealism does not know real, sensuous activity as such. Feuerbach wants sensuous objects, really distinct from the thought what is good and bad, which can vary, while only truth is accorded some kind of objective significance. The categories of ‘truth’, ‘goodness’ and ‘beauty’ were worked out over millennia as objective ideal characteristics of human activity in the world, unfolding (developing) within that activity, and therefore as objective characteristics of the world (of matter). In other words, to put it briefly, when we speak of activity by the standards of truth, goodness and beauty, we are speaking of ‘points of development’ of human sensuousness, cognition and morality, of the development of man’s relation to man and of man’s relation to nature, taken from the side of their ideal dimension. objects, but he does not conceive human activity itself as objective activity.” This remark does more than emphasise the role of idealism in the becoming of the theory of activity—although that role is indeed enormous and important, so much so that Lenin later spoke of the importance of developing societies of “materialist friends of Hegelian dialectics”. The categories of dialectics, understood materialistically as “forms of human activity repeated millions of times” (Lenin), together with their internal contradictoriness and movement, are an indispensable tool for developing the theory of communism precisely because they reflect the essential moments of human activity in general and capture that without which human activity cannot, in principle, take place. But Marx’s thesis also places emphasis on the materiality of human activity.

Let us emphasise once again that free activity is not “leisure” in the modern sense of the word, not idle pastime, but activity by the measure of freedom. Philosophy, in the person of Baruch Spinoza, once defined freedom as recognised necessity. True freedom is defined by acting in accordance with this necessity.

In their free time, in accordance with historical necessity, people develop culture in the broadest sense of the word. Such activity is not something invented by communists in fantastical dreams. It existed in the past and exists now. But it existed, and continues to exist, within very rigid limits and is the privilege of only a few. The overwhelming majority of the members of society do not know free activity. This has been and still is the case, among other reasons, because one needs not only to have the time but also to be developed enough to be capable of carrying out such activity. The conditions for such development have been, and remain, highly limited in the modern society of commodity production.

So, in order to move in the direction of communism, what must be produced is communism itself, that is, free activity, or more precisely, free time and a corresponding level of development of the members of society who could carry it out.

Consequently, distribution here is, first and foremost, the distribution of free time and of opportunities for the assimilation of human culture. The distribution of the products of labour is necessary only insofar as it provides free time. It is in this spirit that the question of the efficiency of social production in the period of transition to communism should be posed. Under capitalism, free time is a cost of production; it is not accounted for, not valued, and therefore often never becomes truly free and is not used for its proper purpose. The purpose of free time is to be used for the continuous development of social relations (or the development of activity, which is essentially the same thing).

That is why, for communism, it is very important that the habit of treating one's own free time, and the activity carried out during it, as social wealth not subject to private squandering should become the common possession of all. In general, the force of habit and the collective formation of habits are of enormous significance here, if only because being free, that is, acting by the measure of freedom, is always very difficult. It means being at the cutting edge of the contradictions through which the development of society takes place. This is not only hard work but also a tragedy, in the sense that one must constantly overcome oneself and negate one's past self. And collective social freedom is a collective social tragedy and its constant resolution. Communist society is not a society of universal laziness and idleness, but a society of universal free activity, which is for people life’s prime want.

Communism is realised through the necessary succession of directly collective forms of human activity aimed at transforming that very activity (at developing people’s life itself). The forms already attained, striving beyond their own limits, give rise to new ones. That is why this succession cannot be pictured as the same metamorphoses endlessly repeating one after another (as we see in the movement of capital). The inner logic of the movement of the directly collective forms of free activity, the necessity of their succession, is determined not merely by the historically attained mode of production and the relations of production corresponding to it, but by freedom as the necessity of matter itself, which unfolds as matter is assimilated.

In its most general form, the law of communism can be formulated as the law of freedom, or the law of the self-determination of human activity: the free activity of the members of society, carried out as the conscious development of the content and forms of that activity, determines the development of society as a whole. Or, in other words, freedom, in realising itself, strives towards self-expansion.

Of course, understanding how this law operates requires constant refinement, drawing on all available knowledge of both nature and society. That is why, under the appropriate conditions, the theory of communism can play a unifying role for all the sciences, pointing to the real connection between them as the connection of human practice itself.

The directly collective forms of free human activity that exist today as “shoots of communism” are not merely evidence of its possibility; they are its real being. It is important to understand that the communist mode of production already exists; the point is only that it should become dominant and extend to the whole of society.

In their movement, the directly collective forms of free human activity, as “shoots of communism”, can unfold into communism as a social structure, that is, as society’s mode of existence. In this sense, one of the most valuable propositions in Lenin’s theoretical legacy is his advice to seek out and nurture the “shoots” of communism wherever they are found.
\end{document}
